\section{Method}
\label{sec:method}

\subsection{Problem Setup}
\label{sec:method_setup}

We study video continuation with a pretrained video diffusion transformer.
Given observed frames, the model generates future frames at 480p using
classifier-free guidance and a fixed denoising schedule. Test-time adaptation
is allowed to optimize only on the observed conditioning window; no future
frames are used during adaptation.

The primary backbone is \backbone~\citep{meituan2025longcatvideo}, a
13.6B-parameter DiT with \nblocks single-stream transformer blocks and hidden
dimension \dimh. The timestep embedding dimension used by the adaLN pathway is
\dimt. Each block maps the timestep representation into shift, scale, and gate
vectors for self-attention and MLP modulation.

\subsection{\method Formulation}
\label{sec:method_formulation}

Let the timestep embedding be $t \in \mathbb{R}^{512}$. \method learns an
additive residual $\delta$ and uses
\begin{equation}
  t' \;=\; t + \delta.
  \label{eq:adasteer-residual}
\end{equation}
Each transformer block then applies its frozen adaLN projection,
\begin{equation}
  t' \;\xrightarrow{\,\mathrm{adaLN}_i\,}\;
  \bigl[
    \gamma_i^{\mathrm{msa}}, \beta_i^{\mathrm{msa}}, \alpha_i^{\mathrm{msa}},\,
    \gamma_i^{\mathrm{mlp}}, \beta_i^{\mathrm{mlp}}, \alpha_i^{\mathrm{mlp}}
  \bigr]
  \;\in\; \mathbb{R}^{6 \times 4096}.
  \label{eq:adaln}
\end{equation}
The same residual $\delta$ is shared across all blocks, but because each block
has its own frozen adaLN projection $\mathrm{adaLN}_i$, the induced modulation
is block-specific. This is the central mechanism of \method and is illustrated
in Figure~\ref{fig:method}.

\begin{figure}[t]
  \centering
  % \includegraphics[width=\linewidth]{method_diagram}
  \rule{0pt}{4cm}\framebox[0.9\linewidth]{\rule{0pt}{4cm}}%
  \caption{\method overview. A single shared residual $\delta$ is learned per
  test clip and added to the global timestep embedding $t$ before each frozen
  per-block adaLN projection consumes it. The frozen adaLN matrices act as a
  learned de-tying mechanism that converts the shared residual into
  block-specific modulation.}
  \label{fig:method}
\end{figure}

\subsection{Structured Weight Tying via Frozen adaLN De-tying}
\label{sec:method_tying}

The trainable surface of \method is a single \dimt-dimensional vector $\delta$,
shared globally across all \nblocks transformer blocks. The induced
perturbation, however, is per-block: each block $i$ consumes $t' = t + \delta$
through its own frozen projection $\mathrm{adaLN}_i$ and produces its own
shift, scale, and gate vectors. We call this arrangement \emph{structured
weight tying}: the trainable degrees of freedom are tied (one shared $\delta$),
but the frozen adaLN matrices supply learned de-tying that lifts that single
update into \nblocks distinct modulation perturbations.

Two consequences follow. First, the adaptation capacity is bounded by the
intrinsic dimensionality of the timestep pathway rather than by the model
width, so the same method scales from smaller DiTs to billion-parameter
backbones without enlarging the trainable surface. Second, because the
modulation projections themselves are unchanged, the method respects the
pretrained relationships between the timestep code and per-block behavior;
$\delta$ can only steer behavior along directions the pretrained backbone
already considers meaningful, which we view as a useful inductive bias for
single-clip TTA where overfitting is the dominant failure mode.

\subsection{Optimization}
\label{sec:method_optim}

\method is optimized with the standard denoising or flow-matching objective on
the conditioning frames. The adaptation loop samples noise levels, corrupts
latent encodings of the observed frames, predicts the denoising target, and
updates only the \method residual $\delta$. The adapted residual is applied
during future-frame generation and discarded after the video is generated.

The default strong configuration is 10 adaptation steps with learning rate
$5\!\times\!10^{-3}$. This configuration was selected on the 200-video
discovery sweep subsets (Section~\ref{sec:exp_discovery}) and re-evaluated at
the 1000-video scale before being reported: moderate learning rates and
5--10 steps preserve PSNR while improving FVD, while higher learning rates or
longer adaptation can overfit.

\subsection{TinyLoRA / SVD Variant}
\label{sec:method_tinylora}

We include a TinyLoRA variant inspired by SVDiff~\citep{han2023svdiff}. The
useful formulation is
\begin{equation}
  y \;=\; Wx \;+\; \frac{\alpha}{r}\, U_r \bigl(\,v \odot (V_r^\top x)\,\bigr).
  \label{eq:tinylora}
\end{equation}
For \backbone attention targets with \nblocks blocks and two modules per block,
rank-2 untied TinyLoRA trains 192 scalars. With full tying across blocks, this
can reduce to 4 scalars. This supports the broader argument that highly
constrained adaptation subspaces are better matched to single-video TTA than
free LoRA factors.
